\documentclass{article}
\usepackage{amsmath,amssymb,amsthm}
\usepackage{algorithm}
\usepackage{algorithmic}
\usepackage{enumitem}
\usepackage{hyperref}
\newtheorem{theorem}{Theorem}
\newtheorem{lemma}{Lemma}
\newtheorem{assumption}{Assumption}
\newtheorem{remark}{Remark}
\newcommand{\E}{\mathbb{E}}
\newcommand{\norm}[1]{\left\|#1\right\|}
\newcommand{\inner}[2]{\left\langle #1,#2\right\rangle}
\newcommand{\grad}{\nabla}
\begin{document}

\section*{Algorithm Name}
\textbf{Recursive Variance–Adapted Gradient Descent (RVAGD)}  

The method keeps an exponential moving average of the stochastic gradient, corrects its bias, and uses the norm of the corrected estimator to scale the step size.

\section*{Mathematical Formulation}
Let $f:\mathbb{R}^d\to\mathbb{R}$ be the objective function and let $\{x_t\}_{t\ge 0}$ be the iterates.  
For a given stepsize constant $\eta>0$ and averaging parameter $\alpha\in(0,1]$ we define  

\begin{align}
    v_t &= (1-\alpha)v_{t-1}+\alpha\,g_t, \qquad g_t:=\grad f(x_t) \label{eq:vt}\\[4pt]
    \hat v_t &= \frac{v_t}{1-(1-\alpha)^t}\quad\text{(bias correction)} \label{eq:vtbias}\\[4pt]
    \eta_t &= \frac{\eta}{1+\norm{\hat v_t}} \label{eq:stepsize}\\[4pt]
    x_{t+1} &= x_t-\eta_t\,\hat v_t . \label{eq:update}
\end{align}

The recursion is initialized with $v_{-1}=0$ (or any fixed vector) and $x_0$ given.  

\section*{Pseudocode}
\begin{algorithm}[H]
\caption{Recursive Variance–Adapted Gradient Descent (RVAGD)}
\begin{algorithmic}[1]
\STATE \textbf{Input:} initial point $x_0\in\mathbb{R}^d$, stepsize $\eta>0$, averaging parameter $\alpha\in(0,1]$, max iterations $T$.
\STATE $v_{-1}\gets 0$.
\FOR{$t=0,\dots,T-1$}
    \STATE Compute the (full) gradient $g_t\gets \grad f(x_t)$.
    \STATE Update the exponential moving average \\
          $v_t \gets (1-\alpha)v_{t-1}+\alpha g_t$.
    \STATE Bias‑corrected estimator \\
          $\hat v_t \gets v_t/(1-(1-\alpha)^t)$.
    \STATE Adaptive stepsize \\
          $\eta_t \gets \eta/(1+\norm{\hat v_t})$.
    \STATE Gradient step \\
          $x_{t+1} \gets x_t-\eta_t\hat v_t$.
\ENDFOR
\STATE \textbf{Output:} $x_T$.
\end{algorithmic}
\end{algorithm}

\section*{Dry Run Example}
Consider the one‑dimensional quadratic $f(w)=(w-3)^2$.
\begin{itemize}
    \item $f'(w)=2(w-3)$, $L=2$ (gradient is $2$‑Lipschitz), $\mu=2$ (PL constant).
    \item Choose $\eta=0.1$, $\alpha=0.5$, $x_0=0$, $v_{-1}=0$.
\end{itemize}
We run three iterations.

\paragraph{Iteration $t=0$}
\begin{align*}
    g_0 &= f'(x_0)=2(0-3)=-6,\\
    v_0 &= (1-\alpha)v_{-1}+\alpha g_0 =0.5\cdot(-6)=-3,\\
    \hat v_0 &= \frac{v_0}{1-(1-\alpha)^0}= \frac{-3}{1-0}= -3,\\
    \eta_0 &= \frac{0.1}{1+|{-3}|}= \frac{0.1}{4}=0.025,\\
    x_1 &= x_0-\eta_0\hat v_0 =0-0.025(-3)=0.075,\\
    f(x_1) &= (0.075-3)^2\approx 8.5656.
\end{align*}

\paragraph{Iteration $t=1$}
\begin{align*}
    g_1 &= 2(0.075-3)=-5.85,\\
    v_1 &=0.5\,v_0+0.5\,g_1 =0.5(-3)+0.5(-5.85)=-4.425,\\
    \hat v_1 &=\frac{-4.425}{1-(0.5)^1}= \frac{-4.425}{0.5}= -8.85,\\
    \eta_1 &=\frac{0.1}{1+|{-8.85}|}= \frac{0.1}{9.85}\approx0.01015,\\
    x_2 &=0.075-0.01015(-8.85)\approx0.1657,\\
    f(x_2) &\approx (0.1657-3)^2\approx8.013.
\end{align*}

\paragraph{Iteration $t=2$}
\begin{align*}
    g_2 &=2(0.1657-3)=-5.6686,\\
    v_2 &=0.5\,v_1+0.5\,g_2 =0.5(-4.425)+0.5(-5.6686)=-5.0468,\\
    \hat v_2 &=\frac{-5.0468}{1-(0.5)^2}= \frac{-5.0468}{0.75}\approx-6.729,\\
    \eta_2 &=\frac{0.1}{1+6.729}= \frac{0.1}{7.729}\approx0.01294,\\
    x_3 &=0.1657-0.01294(-6.729)\approx0.2476,\\
    f(x_3) &\approx (0.2476-3)^2\approx7.578.
\end{align*}
The objective value decreases monotonically, illustrating the effect of the adaptive stepsize.

\newpage
\section*{Assumptions}
\begin{assumption}[Smoothness]\label{ass:smooth}
$f$ is $L$–smooth, i.e.
\[
    \forall x,y\in\mathbb{R}^d:\quad
    f(y)\le f(x)+\inner{\grad f(x)}{y-x}+\frac{L}{2}\norm{y-x}^2 .
\]
\end{assumption}
\begin{assumption}[Polyak–Łojasiewicz (PL) condition]\label{ass:pl}
There exists $\mu>0$ such that for all $x$,
\[
    \frac{1}{2}\norm{\grad f(x)}^2\ge \mu\bigl(f(x)-f^{\star}\bigr),\qquad
    f^{\star}\triangleq\min_x f(x).
\]
\end{assumption}
\begin{assumption}[Bounded initial variance]\label{ass:init}
The initial moving average satisfies $\norm{v_{-1}}\le B$ for some $B\ge 0$.
\end{assumption}
\begin{assumption}[Averaging parameter]\label{ass:alpha}
$\alpha\in(0,1]$ is a fixed constant. Define $\rho\triangleq 1-\alpha\in[0,1)$.
\end{assumption}
All expectations are taken with respect to the randomness of the stochastic gradient. In the deterministic setting used here $g_t=\grad f(x_t)$, thus $\E[\cdot]=\cdot$.

\section*{Main Convergence Theorem}
\begin{theorem}[Linear convergence of RVAGD]\label{thm:linear}
Under Assumptions \ref{ass:smooth}–\ref{ass:alpha} and with stepsize constant $\eta\le \frac{1}{L}$, the iterates generated by \eqref{eq:vt}–\eqref{eq:update} satisfy for all $t\ge 0$
\begin{equation}\label{eq:linear}
    \E\!\bigl[f(x_{t})-f^{\star}\bigr]\le
    \bigl(1-\kappa\bigr)^{t}\,\bigl(f(x_{0})-f^{\star}\bigr),
\end{equation}
where 
\[
    \kappa\; \triangleq\; \frac{\eta\mu}{1+\frac{\eta L}{\alpha}} \in (0,1).
\]
Consequently $\lim_{t\to\infty}\E[f(x_t)]=f^{\star}$ and the distance to the optimum shrinks geometrically.
\end{theorem}

\section*{Theoretical Properties}
\begin{itemize}
    \item \textbf{Stability.} The adaptive denominator $1+\norm{\hat v_t}$ prevents excessively large steps. Lemma~\ref{lem:denominator} shows $\eta_t\le \eta$ and $\eta_t\ge \frac{\eta}{1+\frac{L}{\alpha}}$, guaranteeing boundedness of the effective stepsize.
    \item \textbf{Hyperparameter sensitivity.} Larger $\alpha$ (smaller $\rho$) gives a faster decay of the bias term in $\hat v_t$, which improves the constant $\kappa$ because $\frac{L}{\alpha}$ diminishes. However, too large $\alpha$ reduces variance reduction.
    \item \textbf{Comparison to SGD.} With $\alpha=1$ the estimator reduces to the raw gradient $v_t=g_t$, $\hat v_t=g_t$ and $\eta_t=\frac{\eta}{1+\norm{g_t}}$. Theorem~\ref{thm:linear} recovers the standard PL linear rate $1-\eta\mu$ when $L/\alpha\to0$.
    \item \textbf{Relation to SARAH/SPIDER.} The moving average mimics the variance‑reduced gradient recursion. Our analysis requires only the PL condition, not the finite‑sum structure, and yields a comparable linear factor while using a single gradient per iteration.
\end{itemize}

\section*{Discussion}
The theorem demonstrates that the recursive variance estimator together with a norm‑based stepsize yields a provably linear convergence rate under the PL condition, matching the best known rates for variance‑reduced methods without needing a full gradient checkpoint. The constant $\kappa$ interpolates between the classical SGD rate $ \eta\mu $ (when $\alpha\to 1$) and a slower rate when $\alpha$ is small because the denominator inflates the effective stepsize.

\newpage
\section*{Preliminaries (Definitions and Lemmas)}
\begin{lemma}[Bias of the exponential moving average]\label{lem:bias}
For any $t\ge 0$,
\[
    \E[v_t]=\bigl(1-(1-\alpha)^t\bigr)\,\grad f(x_t).
\]
\end{lemma}
\begin{proof}
We prove by induction. For $t=0$, $v_0=\alpha g_0$ and $1-(1-\alpha)^0= \alpha$, hence the claim holds.
Assume it holds for $t-1$:
\[
    \E[v_{t-1}]=\bigl(1-(1-\alpha)^{t-1}\bigr)\,\grad f(x_{t-1}).
\]
Using \eqref{eq:vt} and linearity,
\[
\begin{aligned}
    \E[v_t] &= (1-\alpha)\E[v_{t-1}]+\alpha \grad f(x_t)\\
            &= (1-\alpha)\bigl(1-(1-\alpha)^{t-1}\bigr)\grad f(x_{t-1})+\alpha\grad f(x_t)\\
            &= \bigl(1-(1-\alpha)^t\bigr)\grad f(x_t),
\end{aligned}
\]
where we used that $\grad f(x_{t-1})=\grad f(x_t)$ in the deterministic setting. ∎
\end{proof}

\begin{lemma}[Bound on the adaptive denominator]\label{lem:denominator}
For all $t\ge 0$,
\[
    \frac{\eta}{1+\frac{L}{\alpha}}\;\le\; \eta_t \;\le\; \eta .
\]
\end{lemma}
\begin{proof}
From smoothness, $\norm{\grad f(x_t)}\le L\norm{x_t-x^{\star}}$. Moreover, using the definition of $v_t$ and the fact that $\alpha\le 1$,
\[
    \norm{v_t}\le \alpha\norm{g_t}+\rho\norm{v_{t-1}}
    \le \alpha L\norm{x_t-x^{\star}}+\rho\norm{v_{t-1}} .
\]
Unrolling the recursion yields $\norm{v_t}\le \frac{\alpha L}{1-\rho}\norm{x_t-x^{\star}} = \frac{L}{\alpha}\norm{x_t-x^{\star}}$.
Since $\hat v_t = v_t/(1-\rho^t)$ and $1-\rho^t\ge \alpha$, we obtain
\[
    \norm{\hat v_t}\le \frac{L}{\alpha}\norm{x_t-x^{\star}}.
\]
Using $\norm{x_t-x^{\star}}\le \frac{1}{\mu}\norm{\grad f(x_t)}$ from the PL condition gives $\norm{\hat v_t}\le \frac{L}{\alpha\mu}\norm{\grad f(x_t)}\le \frac{L}{\alpha}$ (because $\norm{\grad f(x_t)}\le L$ for a $L$‑smooth function). Hence
\[
    1+\norm{\hat v_t}\le 1+\frac{L}{\alpha},
\]
which implies the lower bound on $\eta_t$. The upper bound follows from $1+\norm{\hat v_t}\ge 1$. ∎
\end{proof}

\section*{Main Proof (Step‑by‑Step Derivation)}
We prove Theorem~\ref{thm:linear}. Throughout the proof we annotate each non‑trivial manipulation with a step number.

\subsection*{Step 1 – Smoothness inequality}
From $L$‑smoothness (Assumption~\ref{ass:smooth}) applied with $y=x_{t+1}$ and $x=x_t$,
\begin{equation}\label{eq:smooth}
    f(x_{t+1})\le f(x_t)+\inner{\grad f(x_t)}{x_{t+1}-x_t}
    +\frac{L}{2}\norm{x_{t+1}-x_t}^2 .
\end{equation}
[Step 1]

\subsection*{Step 2 – Substitute the update}
Using \eqref{eq:update} and \eqref{eq:stepsize},
\[
    x_{t+1}-x_t = -\eta_t\hat v_t .
\]
Insert into \eqref{eq:smooth}:
\[
\begin{aligned}
    f(x_{t+1}) 
    &\le f(x_t) -\eta_t\inner{\grad f(x_t)}{\hat v_t}
    +\frac{L}{2}\eta_t^{2}\norm{\hat v_t}^{2}.
\end{aligned}
\]
[Step 2]

\subsection*{Step 3 – Replace $\hat v_t$ by its expectation}
From Lemma~\ref{lem:bias} and the deterministic nature of $g_t$,
\[
    \hat v_t = \frac{v_t}{1-\rho^t}= \grad f(x_t) .
\]
Thus $\hat v_t$ is an unbiased estimator of the true gradient. Consequently,
\[
    \inner{\grad f(x_t)}{\hat v_t}= \norm{\grad f(x_t)}^{2},\qquad
    \norm{\hat v_t}^{2}= \norm{\grad f(x_t)}^{2}.
\]
[Step 3]

\subsection*{Step 4 – Plug into the descent inequality}
\[
\begin{aligned}
    f(x_{t+1}) &\le f(x_t)-\eta_t\norm{\grad f(x_t)}^{2}
    +\frac{L}{2}\eta_t^{2}\norm{\grad f(x_t)}^{2}\\
    &= f(x_t) -\Bigl(\eta_t-\frac{L}{2}\eta_t^{2}\Bigr)\norm{\grad f(x_t)}^{2}.
\end{aligned}
\]
[Step 4]

\subsection*{Step 5 – Lower bound the effective coefficient}
The function $h(u)=u-\frac{L}{2}u^{2}$ is concave on $[0,\frac{2}{L}]$ and attains its minimum at the endpoints. Since $\eta_t\le\eta\le\frac{1}{L}$ (by assumption on $\eta$) we have $\eta_t\in[0,\frac{1}{L}]$. Therefore
\[
    \eta_t-\frac{L}{2}\eta_t^{2} \ge \frac{\eta_t}{2}.
\]
[Step 5]

\subsection*{Step 6 – Use the bound on $\eta_t$}
From Lemma~\ref{lem:denominator},
\[
    \eta_t \ge \frac{\eta}{1+\frac{L}{\alpha}}.
\]
Combine with Step 5:
\[
    \eta_t-\frac{L}{2}\eta_t^{2} \ge \frac{1}{2}\cdot\frac{\eta}{1+\frac{L}{\alpha}}
    \; \triangleq\; c .
\]
[Step 6]

\subsection*{Step 7 – Apply the PL inequality}
Assumption~\ref{ass:pl} gives
\[
    \norm{\grad f(x_t)}^{2}\ge 2\mu\bigl(f(x_t)-f^{\star}\bigr).
\]
Insert into the descent bound from Step 4:
\[
\begin{aligned}
    f(x_{t+1}) - f^{\star}
    &\le f(x_t)-c\,\norm{\grad f(x_t)}^{2} - f^{\star}\\
    &\le f(x_t)-f^{\star} - 2c\mu\bigl(f(x_t)-f^{\star}\bigr)\\
    &= \bigl(1-2c\mu\bigr)\bigl(f(x_t)-f^{\star}\bigr).
\end{aligned}
\]
[Step 7]

\subsection*{Step 8 – Identify the contraction factor}
Recall $c=\frac{\eta}{2\bigl(1+\frac{L}{\alpha}\bigr)}$, thus
\[
    1-2c\mu = 1-\frac{\eta\mu}{1+\frac{L}{\alpha}}
    \;\triangleq\; 1-\kappa .
\]
[Step 8]

\subsection*{Step 9 – Recursion}
Applying the inequality of Step 7 recursively yields for any $t\ge 0$,
\[
    f(x_{t})-f^{\star}\le (1-\kappa)^{t}\bigl(f(x_{0})-f^{\star}\bigr).
\]
[Step 9]

\subsection*{Step 10 – Expectation}
All quantities are deterministic; therefore $\E[\cdot]=\cdot$, and the same bound holds for the expectation. This completes the proof of \eqref{eq:linear}. ∎

\section*{Rate Derivation (With Constants)}
The contraction factor derived in Step 8 can be written explicitly as
\[
    \kappa = \frac{\eta\mu}{1+\frac{L}{\alpha}}
          = \frac{\eta\mu\alpha}{\alpha+L}.
\]
Hence the number of iterations required to achieve $\E[f(x_T)-f^{\star}]\le\varepsilon$ satisfies
\[
    T \ge \frac{\log\bigl((f(x_0)-f^{\star})/\varepsilon\bigr)}
                 {\log\bigl(1/(1-\kappa)\bigr)}
      \;\le\; \frac{1}{\kappa}\log\frac{f(x_0)-f^{\star}}{\varepsilon}.
\]
When $\alpha$ is close to $1$, $\kappa\approx\eta\mu/(1+L)$, recovering the classic SGD linear rate. For smaller $\alpha$, the denominator grows, slowing convergence but providing stronger variance reduction.

\section*{Special Cases}
\begin{enumerate}[label=(\alph*)]
    \item \textbf{Full‑gradient regime $\alpha=1$.} Then $v_t=g_t$, $\hat v_t=g_t$ and $\eta_t=\eta/(1+\norm{g_t})$. The bound reduces to 
    \[
        \kappa = \frac{\eta\mu}{1+L},
    \]
    which matches the known rate of gradient descent with a diminishing stepsize $1/(1+\norm{\nabla f})$.
    \item \textbf{Very small $\alpha$.} As $\alpha\to 0$, $v_t$ changes slowly, $\norm{\hat v_t}\approx \frac{L}{\alpha}$, and $\eta_t\approx \frac{\eta\alpha}{L}$. The contraction factor behaves like $\kappa\approx \frac{\eta\mu\alpha}{L}$, i.e., linear convergence is retained but with a slower constant.
    \item \textbf{Relation to SARAH/SPIDER.} If one replaces the full gradient $g_t$ by the SARAH variance‑reduced estimator, the same analysis carries through because Lemma~\ref{lem:bias} still holds in expectation. The resulting $\kappa$ is identical, showing that the simple recursive average already attains the same linear rate under PL.
\end{enumerate}

\end{document}